\documentclass[10pt]{article}
\usepackage[a4paper, margin=0.72in]{geometry}
\usepackage{titlesec}
\usepackage{booktabs}
\usepackage{tabularx}
\usepackage{xcolor}
\usepackage{enumitem}
\usepackage{tgheros}
\usepackage[T1]{fontenc}
\usepackage{array}
\usepackage{textcomp}
\usepackage[hidelinks]{hyperref}
\usepackage{tikz}
\usetikzlibrary{shapes.rounded, arrows.meta, positioning}
\renewcommand{\familydefault}{\sfdefault}

%% Colour palette
\definecolor{primary}{RGB}{15,52,96}        %% Deep petroleum navy
\definecolor{accent}{RGB}{198,124,0}        %% Oil amber
\definecolor{lightgrey}{gray}{0.85}
\definecolor{darkgrey}{gray}{0.38}

%% Section style
\titleformat{\section}{%
  \color{primary}\bfseries\large\raggedright
}{}{0em}{}[\color{lightgrey}{\titlerule[1.5pt]}\vspace{-5pt}]

\titleformat{\subsection}{%
  \color{accent}\bfseries\normalsize\raggedright
}{}{0em}{}

\setlength{\parindent}{0pt}
\setlength{\parskip}{4pt}
\hypersetup{colorlinks=true, urlcolor=primary}

%% =====================================================
\begin{document}

%% ----- Title block -----
\begin{center}
  \vspace*{0.15cm}
  {\color{primary}\LARGE\bfseries NEWCROSS AND PAN OCEAN GROUP}\\[5pt]
  {\color{primary}\large\bfseries Digital Transformation --- A Strategic Vision}\\[10pt]
  {\color{darkgrey}\normalsize Version 1.0\quad\textbar\quad April 2026\quad\textbar\quad
   Discussion Paper --- For Internal Circulation}
\end{center}

\noindent\colorbox{primary!8}{%
\begin{minipage}{\dimexpr\textwidth-2\fboxsep}
  \vspace{5pt}
  \textbf{\color{primary}Thesis.}\quad
  The Newcross and Pan Ocean Group of Companies already holds the most valuable
  asset in any software business: fifty years of operational data, deep domain expertise,
  and a credibility inside the Nigerian oil and gas industry that no foreign vendor can
  replicate. Digital transformation converts that asset into competitive advantage ---
  first internally, then commercially. The question is not whether to transform.
  The RADA-AMS initiative has already answered that. The question is how to execute
  with discipline, in the right sequence, grounded in Nigerian operational reality.
  \vspace{5pt}
\end{minipage}}

\vspace{10pt}

%% ----------------------------------------------------------
\section{The Group of Companies}
%% ----------------------------------------------------------

The Newcross and Pan Ocean Group of Companies comprises three distinct legal entities
operating with partly shared infrastructure, shared senior leadership, and
complementary asset portfolios across the Nigerian upstream oil and gas sector.

\vspace{4pt}
\small
\begin{tabularx}{\columnwidth}{@{}lXl@{}}
\toprule
\textbf{Entity} & \textbf{Primary Role} & \textbf{Incorporated} \\
\midrule
Pan Ocean Oil Corporation (Nigeria) Ltd
  & Exploration, production, gas processing, and pipeline infrastructure.
    First indigenous company in JV with NNPC; OML-98 producing since 1976.
  & 1973 \\[4pt]
Newcross Petroleum Limited
  & Upstream E\&P operations; corporate strategy and finance.
  & --- \\[4pt]
Newcross Exploration and Production Limited
  & Dedicated exploration and drilling operations, including rig assets
    acquired during the Jonathan-era licensing round.
  & --- \\
\bottomrule
\end{tabularx}
\normalsize

\vspace{4pt}
{\small\color{darkgrey}
\textbf{Leadership.}\enspace
Steven Fadeyi serves as Group Managing Director across all three entities.
Dr.\ Jason Fadeyi has focused on business strategy and finance across the group.
Digital systems designed for this group must support multi-entity operation from
day one --- not as a retrofit. The group is the customer, not any single entity.
}

\vspace{10pt}

%% ----------------------------------------------------------
\section{The Operational Portfolio}
%% ----------------------------------------------------------

\small
\begin{tabularx}{\columnwidth}{@{}lXl@{}}
\toprule
\textbf{Asset} & \textbf{Description} & \textbf{Status} \\
\midrule
OML-98 (Ogharefe field)
  & Original producing block; onshore northern Niger Delta. Operations commenced
    August 1976. Gas utilisation initiative running since 1984.
  & Producing \\[4pt]
OML-147 (Northern Niger Delta)
  & Acquired as OPL-275 in 2007; Production Sharing Contract (PSC) with NNPC
    signed 2009; converted to OML in 2015 following successful exploration work
    programme. Early Production Facility (EPF) at Owa Aladinma: 20,000 bopd +
    70 MMscf/d, 3-phase modularised design.
  & Producing \\[4pt]
Ovade-Ogharefe Gas Plant (OOGP)
  & Gas processing plant; cryogenic LPG facility with 29 storage tanks.
    Cornerstone of West Africa's largest carbon-emission reduction project.
    Phase 2 cryogenic expansion completed 2022.
  & Operating \\[4pt]
Amukpe-Escravos Pipeline (AEPP)
  & 20-inch, 67 km, 160,000 bopd capacity. Built using Continuous Horizontal
    Directional Drilling (CHDD) anti-vandalism technology. Shared-use tariff
    arrangement with neighbouring producers.
  & Operating \\[4pt]
Rig Assets (Newcross E\&P)
  & Drilling rig(s) acquired during Jonathan-era licensing round, held under
    Newcross Exploration and Production Limited.
  & Active \\
\bottomrule
\end{tabularx}
\normalsize

\vspace{4pt}
{\small\color{darkgrey}
\textbf{NNPC relationship.}\enspace
OML-98 was established under an original JV arrangement with NNPC --- that
structure has evolved significantly over the years. OML-147 operates under a live
PSC signed with NNPC in 2009; PSC reporting obligations on OML-147 remain active
and are a direct driver of digital reporting requirements.
}

\vspace{10pt}

%% ----------------------------------------------------------
\section{The Nigerian Oil and Gas Context}
%% ----------------------------------------------------------

{\small
Any serious transformation programme must be grounded in the regulatory and market
environment the group actually operates in. Several recent shifts raise the urgency.
}

\vspace{4pt}
\small
\begin{tabularx}{\columnwidth}{@{}lX@{}}
\toprule
\textbf{Factor} & \textbf{Implication for the Group} \\
\midrule
Petroleum Industry Act 2021 (PIA)
  & Restructured fiscal regime; new royalty and tax framework; Host Community
    Development Trust (HCDT) obligation of 3\% of opex per entity. Tracking and
    reporting HCDT is a legal requirement, not optional. \\[4pt]
NUPRC Regulatory Framework
  & Successor to DPR; expanded mandate, increased reporting cadence and data
    granularity. April 2026: fresh directive mandating standardised templates for
    \textbf{measurement-based methane and GHG reporting} --- OOGP is directly in scope. \\[4pt]
Nigerian Gas Flare Commercialisation Programme (NGFCP)
  & Pan Ocean's gas-flare leadership since 1984 is a credibility asset. Digital
    measurement and reporting of flare volumes is now an active compliance requirement. \\[4pt]
Nigerian Content Act (NCA)
  & Local spend, employment, and training tracking across all three entities and
    both producing blocks. Audit risk is increasing. Currently manual throughout. \\[4pt]
OML-147 PSC Reporting
  & Production, cost recovery, and fiscal data submitted to NNPC under PSC
    terms. Manually generated today from disparate sources. \\[4pt]
Dangote Refinery Effect
  & Domestic crude market dynamics are shifting for the first time in decades.
    Oil marketing, lifting decisions, and buyer allocation now operate in a
    fundamentally different environment. \\
\bottomrule
\end{tabularx}
\normalsize

\vspace{10pt}

%% ----------------------------------------------------------
\section{What Digital Transformation Means Here}
%% ----------------------------------------------------------

Digital transformation is not a technology project. It is the systematic replacement
of human effort and paper with connected systems --- so that data runs the business
instead of people chasing information. It has three stages, and they must happen in
sequence.

\vspace{4pt}
\small
\begin{tabularx}{\columnwidth}{@{}lXX@{}}
\toprule
\textbf{Stage} & \textbf{Definition} & \textbf{Group Reality Today} \\
\midrule
1.\ Digitisation
  & Capture data that exists on paper, in spreadsheets, or in people's heads.
  & 50 years of OML-98 well logs; paper-based HSE records; manual procurement
    tracking across three entities. \\[4pt]
2.\ Digitalisation
  & Connect systems so data flows automatically between them.
  & SAP, SCADA, gas plant instrumentation, and pipeline metering operate in silos.
    Humans extract and transfer data between them manually. \\[4pt]
3.\ Transformation
  & Use connected data to change how decisions are made and how value is created.
  & Real-time cross-entity production visibility; AI-driven predictive maintenance;
    automated regulatory submission. This is the destination. \\
\bottomrule
\end{tabularx}
\normalsize

\vspace{4pt}
\noindent\colorbox{primary!8}{%
\begin{minipage}{\dimexpr\textwidth-2\fboxsep}
  \vspace{5pt}
  \textbf{\color{primary}Where the Group Stands.}\quad
  Stages 1 and 2 are partially complete --- Azure hybrid cloud, SAP, and SharePoint
  are in place. The infrastructure foundation exists. What is missing is the
  integration layer, the data pipelines, and the intelligence layer on top.
  The group is closer to transformation than most Nigerian indigenous operators.
  That is the strategic advantage to accelerate --- not the starting point to rebuild.
  \vspace{5pt}
\end{minipage}}

\vspace{10pt}

%% ----------------------------------------------------------
\section{The Operational Complexity That Demands Change}
%% ----------------------------------------------------------

{\small
The group operates across multiple entities, two producing blocks, a 200 MMscf/d-class
gas plant, a 67 km pipeline, and at least five distinct regulatory reporting frameworks
simultaneously. This complexity is the primary source of digital transformation ROI.
}

\vspace{4pt}
\small
\begin{tabularx}{\columnwidth}{@{}lX@{}}
\toprule
\textbf{Business Function} & \textbf{Current Pain and Opportunity} \\
\midrule
Exploration \& Production
  & Multi-block production data (OML-98 + OML-147) managed across separate
    systems; no unified real-time view; well performance managed manually. \\[4pt]
Gas Plant Operations (OOGP)
  & Cryogenic equipment and LPG storage monitoring is reactive rather than
    predictive. Flare monitoring and GHG reporting is now a live NUPRC
    obligation (April 2026 directive). OOGP with 29 LPG storage tanks is the
    group's highest-consequence asset for unplanned downtime. \\[4pt]
Pipeline Management (AEPP)
  & 67 km corridor through the Niger Delta. Third-party injector metering and
    tariff reconciliation is paper-intensive. Crude theft and encroachment
    detection is reactive. Metering accuracy is a direct revenue issue. \\[4pt]
Oil Marketing
  & Crude lifting schedules, buyer allocation, and PSC equity crude
    reconciliation managed manually. The Dangote refinery's entry into the
    domestic market is adding complexity to a system that was already stretched. \\[4pt]
Finance \& Accounting
  & Multi-currency exposure (Naira + USD); intercompany transactions across
    three entities; NNPC PSC cost recovery accounting; annual budget cycle ---
    all heavily manual with significant reconciliation overhead. \\[4pt]
Regulatory Reporting
  & NUPRC oil and gas returns; PSC reporting to NNPC; HCDT contribution
    tracking; Nigerian Content reports; GHG/methane reporting (fresh April 2026
    directive) --- all generated manually from disparate sources. \\[4pt]
Health, Safety \& Environment (HSE)
  & Permit-to-work systems paper-based. Incident reporting and near-miss
    tracking manual. OOGP is a high-consequence HSE environment by any
    definition (cryogenic processing; LPG storage). \\[4pt]
HR \& Nigerian Content
  & Staff rotation scheduling, training compliance, and local spend
    classification for NCA reporting --- manual across three entities and
    multiple operational sites. \\[4pt]
Procurement \& Supply Chain
  & Remote field operations with significant lead times; FX exposure on
    imported parts; emergency procurement at premium cost; vendor performance
    not systematically tracked. \\[4pt]
IT and Communications Infrastructure
  & LAN, WAN, voice, and remote rig data communications link Lagos HQ,
    Ogharefe, Owa Aladinma, and rig sites. These links are not a background IT
    matter --- they are the physical foundation for every data initiative the
    group will pursue. A degraded WAN link to Ogharefe silences SCADA feeds and
    creates operational blind spots. VSAT links to remote rig sites carry
    WITSML drilling parameter data; a failed link means decisions are made on
    stale information. Connectivity health must be assessed and remediated as a
    prerequisite of the transformation programme. \\
\bottomrule
\end{tabularx}
\normalsize

\vspace{10pt}

%% ----------------------------------------------------------
\section{The Two-Track Approach}
%% ----------------------------------------------------------

\noindent\colorbox{primary!8}{%
\begin{minipage}{\dimexpr\textwidth-2\fboxsep}
  \vspace{5pt}
  \textbf{\color{primary}Core Principle.}\quad
  Build once for the group. Design for the industry from day one. Every internal
  solution should be built with configurability, clean APIs, and multi-entity
  support from the outset --- not retrofitted for broader use after the fact.
  This is the discipline that separates a bespoke IT project from a product.
  \vspace{5pt}
\end{minipage}}

\vspace{6pt}

\small
\begin{tabularx}{\columnwidth}{@{}lXl@{}}
\toprule
\textbf{Track} & \textbf{Purpose} & \textbf{Horizon} \\
\midrule
Track A --- Internal Transformation
  & Efficiency, automation, and AI-driven operations across all group entities.
    Build the platform, prove the value, generate the data.
  & Months 0--24 \\[4pt]
Track B --- Commercialisation
  & Package successful internal solutions as configurable products for the
    Nigerian oil and gas market. The group's first paying customer is itself.
    The group's 50-year operating history is the most powerful reference site
    in Nigeria.
  & Months 12--36 \\
\bottomrule
\end{tabularx}
\normalsize

\vspace{10pt}

%% ----------------------------------------------------------
\section{Track A: Phase Roadmap}
%% ----------------------------------------------------------

\subsection*{Phase 1 --- Quick Wins (Months 0--6)}

{\small
Focus on high-visibility, lower-complexity improvements that establish data
pipelines and generate internal credibility before AI is introduced.}

\vspace{4pt}
\small
\begin{tabularx}{\columnwidth}{@{}lXl@{}}
\toprule
\textbf{Initiative} & \textbf{What It Does} & \textbf{Primary Value} \\
\midrule
Communications and Connectivity Audit
  & Assess WAN, LAN, VSAT, and rig data link capacity, reliability, and
    redundancy across all sites (Lagos HQ, Ogharefe, Owa Aladinma, rig sites).
    Identify gaps that would prevent data pipeline operation. Runs in parallel
    with the data audit --- not after it.
  & Infrastructure prerequisite; prevents data blind spots \\[3pt]
AEPP Tariff Reconciliation System
  & Automate metering, injection tracking, and tariff billing for third-party
    injectors on the Amukpe-Escravos Pipeline.
  & Revenue assurance; removes manual reconciliation errors \\[3pt]
Regulatory Reporting Automation
  & Auto-generate NUPRC oil and gas returns from SAP + operational data for
    OML-98 and OML-147, including new measurement-based GHG/methane templates
    (NUPRC April 2026 directive).
  & Legal compliance; hours $\to$ minutes per submission \\[3pt]
HCDT Contribution Tracker
  & Track and report Host Community Development Trust obligations (3\% of
    opex under PIA 2021) across all three group entities with full audit trail.
  & Legal compliance; audit readiness \\[3pt]
HSE Digital Permit-to-Work
  & Replace paper-based permit-to-work system. Critical priority at OOGP
    given cryogenic operations and LPG storage environment.
  & Safety compliance; liability reduction \\[3pt]
Unified Operations Dashboard (MVP)
  & Single real-time view of OML-98, OML-147, OOGP, and AEPP production
    and status across all entities.
  & Faster management decisions; cross-entity visibility \\[3pt]
Workflow Automation
  & Digitise departmental approval flows currently running through email
    and SharePoint across all three entities.
  & Cycle time reduction; audit trail \\
\bottomrule
\end{tabularx}
\normalsize

\vspace{8pt}
\subsection*{Phase 2 --- AI-Driven Operations (Months 6--18)}

{\small
Build on Phase 1 data and internal credibility to introduce machine learning
where the data and infrastructure are now ready.}

\vspace{4pt}
\small
\begin{tabularx}{\columnwidth}{@{}lX@{}}
\toprule
\textbf{Initiative} & \textbf{Description} \\
\midrule
Predictive Maintenance --- OOGP (Highest Priority Asset)
  & Instrument compressors, cryogenic processing units, and LPG storage tanks with
    vibration, temperature, and pressure sensors. Train time-series anomaly detection
    models on sensor readings and SAP maintenance history. Output: failure probability
    scores, recommended maintenance windows, escalation alerts. One avoided plant
    shutdown justifies the full investment. \\[4pt]
AEPP Pipeline Security Intelligence
  & Flow anomaly detection across all injection and delivery points --- flag pressure
    drops and flow irregularities indicative of theft taps or structural issues.
    Satellite imagery change detection along 67 km Right of Way corridor. GPS-
    coordinated alert system for field response teams. \\[4pt]
Oil Marketing Intelligence
  & Lifting schedule management, buyer allocation tracking, and PSC equity crude
    reconciliation with NNPC. Market pricing dashboard calibrated for the evolving
    domestic crude landscape post-Dangote. \\[4pt]
Nigerian Content \& HR Compliance Tracker
  & Classify spend and employment by NCA category automatically across all three
    entities. Auto-generate NCA and HCDT reports. Flag compliance gaps ahead of
    audit cycles. \\[4pt]
OML-147 PSC Cost Recovery Automation
  & Automate data extraction, calculation, and formatting for NNPC PSC cost
    recovery submissions. Replaces multi-day manual process with an auditable
    pipeline. \\[4pt]
GHG Monitoring Pipeline --- OOGP
  & Sensor-to-report pipeline specifically for the new NUPRC measurement-based
    methane and GHG reporting requirement. Continuous monitoring; auto-generation
    of standardised NUPRC templates. \\
\bottomrule
\end{tabularx}
\normalsize

\vspace{8pt}
\subsection*{Phase 3 --- Foundation for Scale (Months 12--24)}

\small
\begin{itemize}[leftmargin=1.5em, itemsep=2pt]
  \item Cloud-native architecture standards established for all new development across the group
  \item API-first design so internal systems can be repackaged as configurable external products
  \item Data lake consolidating operational, financial, and compliance data across all three entities
  \item Security and access control framework appropriate for multi-entity, future multi-tenant operation
  \item Document Intelligence: OCR and semantic search across 50 years of OML-98 well logs,
        drilling reports, and geological surveys --- unlock the group's institutional memory
  \item OML-147 production optimisation (data collection phase completes; ML modelling begins)
\end{itemize}
\normalsize

\vspace{10pt}

%% ----------------------------------------------------------
\section{Track B: Commercialisation}
%% ----------------------------------------------------------

{\small
The primary market is the Nigerian oil and gas ecosystem. Every product below is
built first for the group's own operations. The group's 50-year operating history
and cross-entity reference deployment is the most powerful sales tool any software
vendor in this market could hold.
}

\vspace{4pt}
\small
\begin{tabularx}{\columnwidth}{@{}lXl@{}}
\toprule
\textbf{Product} & \textbf{Problem Solved} & \textbf{Commercial Rationale} \\
\midrule
\textbf{RegulatoryReport}
  & Auto-generation of NUPRC/DPR statutory returns, including measurement-based
    GHG/methane templates and OML-147 PSC data submissions.
  & Mandatory for every operator; frequency and complexity are increasing \\[4pt]
\textbf{NigerianContent.ai}
  & NCA and HCDT compliance tracking, calculation, and audit trail across
    multi-entity operations.
  & Mandatory under NCA and PIA 2021; currently all manual industry-wide \\[4pt]
\textbf{FieldMaintain}
  & Predictive maintenance SaaS for oil field and gas plant equipment.
    OOGP is the reference deployment.
  & Every Nigerian operator with process plant has this problem \\[4pt]
\textbf{FieldWatch}
  & Pipeline integrity monitoring, encroachment detection, and third-party
    metering reconciliation. AEPP is the 67 km reference deployment.
  & Crude theft costs the industry billions annually \\[4pt]
\textbf{HSETrack}
  & Digital permit-to-work, incident reporting, and near-miss tracking
    designed for Nigerian site conditions.
  & Foreign HSE software does not accommodate Nigerian operational reality \\[4pt]
\textbf{WellOptimise}
  & Production optimisation recommendations from historical well data ---
    choke settings, injection schedules, pressure management.
  & Marginal improvements across two OML blocks translate directly to revenue \\
\bottomrule
\end{tabularx}
\normalsize

\vspace{4pt}
\small
\begin{tabularx}{\columnwidth}{@{}lX@{}}
\toprule
\textbf{Target Customer Type} & \textbf{Why They Would Adopt} \\
\midrule
Indigenous operators (Seplat, Aiteo, Shoreline, etc.)
  & Same operational problems; trust Pan Ocean's credibility over foreign vendors \\
IOC Nigerian subsidiaries
  & Need locally-compliant solutions; frustrated by foreign software built for
    Western regulatory environments \\
Gas companies and midstream operators
  & Same infrastructure and compliance challenges \\
NNPC subsidiaries
  & Regulatory reporting and local content requirements at scale \\
\bottomrule
\end{tabularx}
\normalsize

\vspace{10pt}

%% ----------------------------------------------------------
\section{The Compliance Imperative}
%% ----------------------------------------------------------

{\small
Several initiatives in the Phase 1 roadmap are not discretionary choices. They are
responses to legal obligations already in force or recently activated. This reframes
the investment from ``efficiency project'' to ``compliance infrastructure.''
}

\vspace{6pt}

\noindent\colorbox{primary!8}{%
\begin{minipage}{\dimexpr\textwidth-2\fboxsep}
  \vspace{5pt}
  \textbf{\color{primary}April 2026 --- NUPRC Methane and GHG Directive.}\quad
  NUPRC issued a fresh directive mandating standardised templates and the transition
  to \textbf{measurement-based methane and GHG reporting} across all upstream
  operations. OOGP is directly in scope. Manual compliance with this directive is
  operationally untenable. An automated sensor-to-report pipeline is the only
  sustainable response.
  \vspace{5pt}
\end{minipage}}

\vspace{6pt}
\small
\begin{tabularx}{\columnwidth}{@{}lXl@{}}
\toprule
\textbf{Obligation} & \textbf{Scope} & \textbf{Status} \\
\midrule
HCDT Contribution Tracking
  & 3\% of opex per PIA 2021; applies across all three group entities
  & In force since PIA enactment \\[3pt]
NUPRC Methane / GHG Reporting
  & Measurement-based; standardised templates; OOGP directly in scope
  & Fresh directive, April 2026 \\[3pt]
Nigerian Content Act Reporting
  & Local spend, employment, and training across all entities and blocks
  & Ongoing; audit exposure growing \\[3pt]
OML-147 PSC Reporting to NNPC
  & Production, cost recovery, and fiscal data under live PSC terms
  & Ongoing; live obligation \\[3pt]
Gas Flare Reporting (NGFCP)
  & Flare volume measurement and programme participation
  & Ongoing \\[3pt]
NDPR Compliance
  & Nigeria Data Protection Regulation applies to employee data and
    some operational datasets. Any digital system must be designed
    for NDPR compliance from the outset.
  & Ongoing \\
\bottomrule
\end{tabularx}
\normalsize

\vspace{10pt}

%% ----------------------------------------------------------
\section{Team Architecture}
%% ----------------------------------------------------------

\begin{center}
\begin{tikzpicture}[
  node distance=0.8cm and 2.2cm,
  box/.style={draw, rounded corners=5pt, minimum width=3.2cm,
              minimum height=0.9cm, align=center, font=\small},
  arr/.style={-{Stealth[length=5pt]}, thick, color=primary},
  lbl/.style={font=\scriptsize, color=darkgrey},
]

\node[box, fill=primary!12, draw=primary!70, minimum width=7cm] (head)
  {\textbf{Head of Digital Transformation}\\
   {\scriptsize Strategy \textbullet\ Architecture \textbullet\ Commercial track}};

\node[box, fill=accent!10, draw=accent!60, below left=0.9cm and 2.8cm of head] (ba)
  {\textbf{Business Analyst}\\{\scriptsize Domain \textbullet\ Requirements}};

\node[box, fill=primary!8, draw=primary!50, below=0.9cm of head] (tl)
  {\textbf{Technical Lead}\\{\scriptsize Architecture \textbullet\ Standards}};

\node[box, fill=accent!10, draw=accent!60, below right=0.9cm and 2.8cm of head] (pm)
  {\textbf{Product Manager}\\{\scriptsize Track B \textbullet\ Market}};

\node[box, fill=primary!5, draw=primary!30, below left=0.9cm and 1.5cm of tl, minimum width=3.0cm] (be)
  {\textbf{Backend (2--3)}\\{\scriptsize Python \textbullet\ Java \textbullet\ APIs}};

\node[box, fill=primary!5, draw=primary!30, below=0.9cm of tl, minimum width=3.0cm] (de)
  {\textbf{Data Engineer}\\{\scriptsize Pipelines \textbullet\ ML}};

\node[box, fill=primary!5, draw=primary!30, below right=0.9cm and 1.5cm of tl, minimum width=3.0cm] (fe)
  {\textbf{Frontend (1--2)}\\{\scriptsize React \textbullet\ React Native}};

\node[box, fill=accent!5, draw=accent!30, below=2.3cm of tl, minimum width=3.0cm] (devops)
  {\textbf{DevOps / Cloud}\\{\scriptsize Azure \textbullet\ Kubernetes}};

\draw[arr] (head) -- (ba);
\draw[arr] (head) -- (tl);
\draw[arr] (head) -- (pm);
\draw[arr] (tl) -- (be);
\draw[arr] (tl) -- (de);
\draw[arr] (tl) -- (fe);
\draw[arr] (tl) -- (devops);

\end{tikzpicture}
\end{center}

\vspace{4pt}
{\small\color{darkgrey}
\textbf{Hiring principle.}\enspace Mix internal hires (domain knowledge, institutional
trust) with external engineers (modern practices, cloud-native experience). Prioritise
diaspora engineers: shared cultural context combined with global technical exposure.
Invest in DevOps early --- this is what separates products from projects.
The Business Analyst role requires genuine operational knowledge of Pan Ocean's
processes; this is the hardest role to fill externally and the most important to
get right.
}

\vspace{10pt}

%% ----------------------------------------------------------
\section{Common Failure Modes}
%% ----------------------------------------------------------

\small
\begin{tabularx}{\columnwidth}{@{}lXX@{}}
\toprule
\textbf{Risk} & \textbf{What Happens} & \textbf{Mitigation} \\
\midrule
Building for presentations
  & Impressive demos that don't survive contact with real operations
  & Ship working software early; iterate from actual field usage \\[3pt]
Ignoring infrastructure constraints
  & Solutions requiring constant internet at Ogharefe, Owa Aladinma,
    or along the AEPP corridor fail immediately
  & Offline-first architecture from day one; non-negotiable \\[3pt]
Over-relying on foreign consultants
  & They don't understand NUPRC, the Niger Delta, Naira FX, or
    Nigerian site conditions; solutions fail on contact with reality
  & Build internal capability first; use external for specific
    skills only, never for strategy \\[3pt]
Single-entity design
  & Systems designed only for Pan Ocean that cannot accommodate
    Newcross entities; commercialisation becomes impossible
  & Multi-entity architecture from day one; the group is the customer \\[3pt]
Building bespoke instead of product
  & Every solution is unique; nothing is reusable across entities
    or sellable externally
  & API-first; configuration over customisation from the start \\[3pt]
Starting too broad
  & Ten initiatives in parallel; none prove out; credibility lost
  & Phase approach; fully complete one initiative before scaling \\[3pt]
Ignoring data quality
  & AI trained on bad data produces bad and potentially dangerous decisions
  & Data audit and governance programme before any ML project begins \\
\bottomrule
\end{tabularx}
\normalsize

\vspace{10pt}

%% ----------------------------------------------------------
\section{Guiding Principles}
%% ----------------------------------------------------------

\small
\begin{tabularx}{\columnwidth}{@{}lX@{}}
\toprule
\textbf{Principle} & \textbf{What It Means in Practice} \\
\midrule
Nigerian problems first
  & Build for the NUPRC, the Niger Delta, the Naira, and the generator.
    Not for Silicon Valley assumptions about infrastructure and connectivity. \\[3pt]
Offline-first
  & Core functions must operate without reliable internet at every
    operational site (Lagos HQ, Ogharefe, Owa Aladinma, AEPP corridor). \\[3pt]
Explainable AI
  & Operations managers must understand why the system is recommending
    something. Black-box AI will not be adopted by field teams. \\[3pt]
Data before models
  & A data audit and quality programme precedes any machine learning
    initiative. Dirty data produces dangerous recommendations. \\[3pt]
Locally hosted where sensitive
  & Well data, production history, and commercial data remain on-premise
    or within the group's Azure tenant. Not on third-party cloud APIs. \\[3pt]
Price in Naira
  & Minimise foreign exchange exposure for internal operations and for
    future customers. This is a genuine competitive advantage in the market. \\[3pt]
RADA-AMS is the mandate
  & The group publicly committed to energy and innovation leadership through
    RADA-AMS in October 2025. This work is the execution engine for that
    commitment --- not a new idea, but the operationalisation of an
    existing strategic direction. \\
\bottomrule
\end{tabularx}
\normalsize

\vspace{10pt}

%% ----- Footer -----
\vspace{6pt}
{\small\color{darkgrey}
\textbf{Newcross and Pan Ocean Group of Companies}\enspace|\enspace
Strategic Vision --- Version 1.0\enspace|\enspace April 2026\enspace|\enspace
Discussion Paper\\
Companion document: \textit{AI and Data Roadmap --- Newcross and Pan Ocean Group
(April 2026)}.\\
Prepared drawing on direct operational knowledge of group infrastructure,
NUPRC regulatory updates (April 2026), PIA 2021 obligations, and public
operational data.
}

\end{document}
